\begin{center}
\LARGE
THE BUCHBERGER ALGORITHM TO STUDY THE A
 PRIORI IDENTIFIABILITY OF BIOLOGICAL SYSTEMS
                           MODELS
\end{center}

\begin{center}
\Large
Maria Pia Saccomani, Stefania Audoly, Claudio Cobelli, Leontina D'angi�
\end{center}

\noindent
Date:  July 19th (Friday)\\
Time:  11:00-11:30\\

\begin{center}
\large
Abstract
\end{center}

In this paper we present an application of the Buchberger algorithm to the solution of a crucial step
in the identification of compartmental models of biological systems from input-output
experiments: the a priori identifiability problem. 

Linear compartmental models are a class of dynamic models based on mass conservation principles
which are widely used in biology and medicine to study the kinetics of substances [4,5].
Compartmental models are mathematically described by ordinary differential equations, where the
state variables usually are the compartment masses and the parameters are the transfer rate
constants among them. The input-output experiment usually consists in injecting some material
into one (or more) compartment and measuring a linear function of some of the state variables. 

The a priori structural identifiability of linear compartmental models consists in testing if, under
the ideal conditions of noise-free observations and error-free model structure (a priori) and
independently of the particular values of the parameters (structurally), the unknown parameters of
the model can or can not be uniquelly (globally) estimated from the designed experiment [1-4]. A
priori identifiability is clearly a necessary prerequisite for well posedness of parameter estimation.
It is also crucial in qualitative experiment design [10], that is selecting the input-output
configuration necessary to ensure unique identifiability of the unknown parameters.Various
methods to test a priori structural identifiability have been proposed, including the transfer
function method [1,2], the transfer function topological method [6], the modal matrix method
[2,7], and the similarity transformation method [2,3]. However, whatever is the method used there
is the need to solve a system of nonlinear algebraic equations (the identifiability equations) which
fast increases in number of unknowns, terms and nonlinearity degree as the model connectivity
and/or the number of compartments increase. While some identifiability results are available on
specific classes of models, there is no algorithm to test global identifiability of compartmental
models of general structure. Recently, a differential algebra algorithm, implemented in Maple, has
been proposed [8] for low dimension nonlinear models: however its applicability to linear
compartmental models is severely limited by computability bounds (up to three compartments). 

The aim of our work has been to develop a new algorithm for checking a priori global
identifiability of linear compartmental models of general structure by combining one classical
method with computer algebra methods, in particular with the Gr\"obner basis. In principle all the
methods discussed above could be used to generate the identifiability equations, but we have
chosen the one which makes the algorithm to calculate the Gr\"obner basis successful for the largest
class of models, i.e. the one which reduces the complexity of the identifiability equations in terms
of number of unknowns, number of terms, and nonlinearity degree. The algorithm is essentially a
two stage one. First, the transfer function topological method [6] is used which is based on the
transfer function and on the topological properties of the graph associated to the model. In
particular, the method defines as unknowns in the identifiability equations not the model
parameters (transfer rate constants), but some topological functions of them, i.e. the cycles and
paths of the compartmental graph. The mapping of the compartmental parameter space into that of
cycles and paths, considerably reduces the number of terms and the number of the high degree
terms. Then, this new system of equations is handled by using the Buchberger algorithm [9] which
calculates the Gr\"obner basis. By substituting the model parameters to the cycles and paths in this
basis and applying a second time the Buchberger algorithm, it is possible to determine, for each
model parameter, if there is one, more than one (and how many), or an infinite number of
solutions. The domain of validity of the algorithm is essentially due to the limits of applicability
of the Buchberger algorithm in solving the identifiability equations and is thus difficult to be
established rigorously. In particular, the Buchberger algorithm has limitations on the number of
simultaneous equations to solve, number of unknowns, number of terms and on the nonlinearity
degree. Note that the choice of the order relation plays an important role in the computational
complexity of the algorithm. We thus checked the algorithm in testing a priori identifiability of
the models present in the biomedical literature. From our experience, we can state that, for the
first time, it is possible to test automatically a priori structural global identifiability of linear
compartmental models characterized by the most general muti input-multi output experimental
configuration and by a structure of relatively large connectivity and dimension, i.e. up to twelve
compartments.The algorithm described has been implemented in the software GLOBI (GLOBI for
GLOBal Identifiability) which is written in PASCAL 6.0 and REDUCE 3.5 [11], where a Gr\"obner
basis package is available with the Buchberger algorithm. In particular, for the user graphical
interface, GLOBI utilizes DELPHI [12], an object-oriented program allowing WINDOWS
application to PASCAL programs. \\

{\Large REFERENCES}

\begin{enumerate}

\item
Cobelli C. and J.J. DiStefano III. Parameter and structural identifiability concepts and
    ambiguities: a critical review and analysis. Am. J. Physiol. 239:R7-R24, 1980. 


\item
Godfrey K.R. and J.J. DiStefano III. Identifiability of model parameters. In Identifiability of
    Parametric Models, E. Walter, Ed. Oxford: Pergamon Press, chapt. 1, pp. 1-20, 1987. 


\item
Walter E. Identifiability of State Space Models. Springer, Berlin, 1982. 


\item
Carson E.R., Cobelli C. and L. Finkelstein.The Mathematical Modeling of Metabolic and
    Endocrine Systems. New York: Wiley, 1983. 


\item
Jacquez J.A. Compartmental Analysis in Biology and Medicine. Ann Arbor: The University
    of Michigan Press, 1985. 


\item
Audoly S. and L. D'Angio'. On the identifiability of linear compartmental systems: a
    revisited transfer function approach based on topological properties. Math. Biosci.
    66:201-228, 1983. 


\item
Norton J.P. Normal mode identifiability analysis of linear compartmental systems in linear
    stages. Math. Biosci. 50:95-115, 1980. 


\item
Glad T. and L. Liung. On global identifiability for arbitrary model parametrisation.
    Automatica, 30.2: 265-276, 1994. 


\item
Buchberger B. An algorithmical criterion for the solvability of algebraic system of equation.
    Aequationes Mathematicae, 4, 3:45-50, 1988. 


\item
Saccomani M.P. and C. Cobelli. Qualitative experiment design in physiological system
    identification. IEEE Control Systems Magazine, vol. 12, n. 6, 18-23, 1992. 


\item
Mueller H.M. Algebraic Computing with REDUCE. M.A.H. Mac Callum \& F.J. Wright,
    Clarendom Press, Oxford, 1991. 


\item
Matcho J., D.R. Faulkner, et al. Using DELPHI. Publisher R. Elgey, QUE Corporation,
    1995. 

\end{enumerate}



